
\section{Methodology}
\label{sec:methodology}

\subsection{Problem Formulation}

We formulate unsafe-behavior monitoring as a binary classification
task. For each episode $i$, the monitor receives an input representation
$x_i$ and produces a score $s_i$ indicating the likelihood that the
episode is unsafe. Given a threshold $\tau$, the predicted label is

\begin{equation}
\hat{y}_i =
\begin{cases}
1, & s_i \geq \tau,\\
0, & s_i < \tau,
\end{cases}
\label{eq:prediction}
\end{equation}

where $y_i \in \{0,1\}$ is the reference label and $\hat{y}_i$ is the
monitor prediction.

We compare two input conditions:

\begin{align}
x_i^{\mathrm{out}} &=
(\mathrm{prompt}_i,\mathrm{output}_i),\\
x_i^{\mathrm{cot}} &=
(\mathrm{prompt}_i,\mathrm{cot}_i,\mathrm{output}_i).
\end{align}

The primary detection metric is recall at a target false-positive rate
(FPR) of 5\%. Recall is defined as

\begin{equation}
\mathrm{Recall} = \frac{TP}{TP+FN},
\label{eq:recall}
\end{equation}

and the false-positive rate is

\begin{equation}
\mathrm{FPR} = \frac{FP}{FP+TN}.
\label{eq:fpr}
\end{equation}

\subsection{Monitoring Conditions}

In the output-only condition, the monitor can inspect the prompt and
final model output. In the CoT-plus-output condition, it can inspect
the prompt, the available reasoning trace, and the final output.

The repository configuration identifies these conditions as
\texttt{output\_only} and \texttt{cot\_plus\_output}, respectively.

\subsection{Unsafe-Behavior Categories}

The benchmark groups unsafe episodes into four categories:

\begin{itemize}
    \item \textbf{Reward hacking:} exploiting a reward or evaluation
    mechanism rather than satisfying its intended objective.

    \item \textbf{Deception:} misleading a user, evaluator, or system
    about relevant behavior or information.

    \item \textbf{Sabotage:} intentionally interfering with a process,
    system, or objective.

    \item \textbf{Harmful compliance:} complying with a request that
    would facilitate harmful activity.
\end{itemize}

\subsection{Monitor and Experimental Configuration}

The current repository configuration specifies a keyword-based
monitoring backend. The configuration also specifies a random seed
of 42, a target FPR of 0.05, and a calibration fraction of 0.5.

The episode data are loaded from the configured JSONL dataset.
These values should be checked against the exact configuration and
dataset version used for the final experiment before submission.

\subsection{Evaluation Procedure}

The experiment separates calibration and evaluation data according
to the configured split. A decision threshold is selected using the
calibration data to target the specified FPR. The monitor is then
evaluated on the held-out evaluation data.

The two monitoring conditions are compared using the same
evaluation protocol.

Operational performance is measured separately from detection
quality. Latency measurements are collected under the configured
request rates. Latency summaries and plots should be reported
independently from classification metrics; a larger number of
latency observations does not imply a larger number of labeled
detection episodes.

\subsection{Reproducibility}

To reproduce the experiment, a researcher should record the
repository revision, configuration file, dataset version, random
seed, monitor backend, visible input fields, calibration split,
target FPR, software dependencies, and execution commands.

Any change to these settings should be documented because it may
affect the results.
